\section{Full Human Rubric with Examples}
\label{app:rubric}

This appendix provides the complete, operationalized rubric employed by trained human raters throughout the evaluation pipeline described in \cref{sec:scoring_rubrics,sec:benchmark_harness}. The rubric instruments two complementary evaluation axes aligned with the study's core criteria (\cref{sec:evaluation_criteria}): the \textbf{Socratic Quality Rubric} assesses pedagogical effectiveness per \cref{subsec:pedagogical_quality}, while the \textbf{Context Fidelity Rubric} quantifies contextual adaptability per \cref{subsec:contextual_adaptability}. Both rubrics were applied to de-identified transcripts sampled from the benchmark harness outputs (\cref{subsubsec:data_persistence}) under the blinded and adjudicated protocol detailed in \cref{subsec:sampling_blinding,subsubsec:evaluation_protocol}.

Each dimension is operationalized with explicit scoring anchors and representative examples drawn from pilot ratings to ensure inter-rater reliability (IRR) met the target thresholds ($\kappa \geq 0.6$, ICC(2,$k$) $\geq 0.7$; see \cref{subsec:inter_rater_reliability}). The examples provided below illustrate score levels 1 (poor), 3 (adequate), and 5 (excellent), enabling consistent application across the full corpus of 120 matched prompts analyzed in \cref{sec:statistical_analysis,sec:results}. Raters received these scoring definitions during initial training and calibration sessions, with refresher exercises conducted before each evaluation wave to maintain consistency. All ratings were logged with timestamps, rater identifiers (anonymized), and adjudication flags per the provenance chain in \cref{subsec:data_sources}.

\paragraph{Evaluation context and sampling procedure.} The rubric was applied to a stratified sample drawn from the benchmark corpus generated by the harness (\cref{sec:benchmark_harness}). Each of the five candidate models—GPT-4o mini, Claude 3 Sonnet, Gemini 2.5 Flash, Command R 08-2024, and Llama 4 Maverick 17B (\cref{sec:candidate_models})—produced responses to 120 prompts spanning twelve curricular categories (\cref{subsec:datasets_prompt_bank}). From this pool, we sampled a balanced subset ensuring at least $n=10$ prompts per category per model, with fixed seed reproducibility recorded in \path{data/statistical_summary.json}. Transcripts presented to raters were shuffled and stripped of model identifiers, cost metadata, and latency telemetry to prevent anchoring bias (\cref{subsec:sampling_blinding}).

\paragraph{Inter-rater reliability and adjudication protocol.} Each transcript received independent ratings from two trained evaluators. Disagreements exceeding pre-specified thresholds—absolute difference $\geq 2$ on any single dimension or mean difference $\geq 1.5$ across dimensions—triggered third-rater adjudication per \cref{subsubsec:evaluation_protocol}. The final adjudicated scores were merged back to the harness outputs via stable identifiers (\verb|prompt_id|, model, run timestamp) as described in \cref{subsec:data_sources}, enabling reproducible aggregation in \cref{subsec:aggregation_weighting} and statistical analysis in \cref{sec:statistical_analysis}. IRR estimates (\cref{subsec:inter_rater_reliability}) confirmed adequate agreement: weighted Cohen's $\kappa=0.469$ (Socratic) and $\kappa=0.498$ (Context), with ICC(2,$k$)=0.73 and 0.75 respectively, meeting acceptance criteria.

\paragraph{Relationship to automated metrics.} While the rubric captures nuanced pedagogical and contextual judgments, automated proxies (\cref{subsec:automated_metrics}) were computed in parallel to triangulate findings. Metrics such as question ratio, keyword coverage, numeric fidelity, and PII flag rate supplemented human scores but did not substitute for them. Discrepancies between automated proxies and human ratings (e.g., high question ratio paired with low Socratic quality score) flagged superficial questioning patterns and were investigated qualitatively in \cref{subsec:qualitative_analysis}.

\import{appendices/A_Full_Human_Rubric/}{A_1_Socratic_Quality.tex}
\import{appendices/A_Full_Human_Rubric/}{A_2_Context_Fidelity.tex}
